\section{Tổng quan hệ thống}
\label{section:components}

RoomHub lấy ứng dụng di động làm trung tâm để giảng viên và cán bộ quản lý cùng thực hiện quy trình mượn phòng. Bên cạnh ứng dụng là dịch vụ đồng bộ dữ liệu và các công cụ hỗ trợ kiểm tra giao tiếp với khóa thông minh. Vai trò của từng thành phần được tóm tắt trong Bảng~\ref{tab:components}; vị trí các phần chính trong kho mã nguồn được trình bày tại Phụ lục~\ref{appendix:components}.

\begin{table}[H]
\centering
\caption{Các thành phần chính của hệ thống RoomHub}
\label{tab:components}
\small
\begin{tabularx}{\textwidth}{>{\raggedright\arraybackslash}p{3.4cm} >{\raggedright\arraybackslash}p{2.7cm} X}
\toprule
\rowcolor{RoomHubLight}
\textbf{Thành phần} & \textbf{Công nghệ} & \textbf{Vai trò} \\
\midrule
Ứng dụng di động & React Native, TypeScript, Kotlin & Cung cấp giao diện theo vai trò, xử lý nghiệp vụ, lưu dữ liệu cục bộ và giao tiếp với OneIoT. \\
Dịch vụ đồng bộ & Java, Spring Boot, H2 & Cung cấp giao diện dữ liệu cho tài khoản, phòng, lịch đặt, bảo trì, quyền truy cập và thông báo. \\
Mô phỏng khóa thông minh & Python & Tạo môi trường kiểm tra bản tin cấp mật khẩu và sự kiện truy cập khi chưa dùng khóa vật lý. \\
Công cụ giao tiếp & Python, MQTT & Làm mẫu tham chiếu cho cơ chế kết nối và định dạng bản tin trao đổi với OneIoT. \\
\bottomrule
\end{tabularx}
\end{table}

Hình~\ref{fig:architecture} thể hiện quan hệ giữa các thành phần. Trong ứng dụng, màn hình tiếp nhận thao tác của người dùng; lớp nghiệp vụ kiểm tra quyền và quy tắc đặt phòng; lớp lưu trữ giữ dữ liệu cục bộ hoặc trao đổi với dịch vụ đồng bộ theo chế độ hoạt động.

\begin{figure}[H]
\centering
\includegraphics[width=0.78\textwidth]{kien_truc.pdf}
\caption{Kiến trúc tổng thể của hệ thống RoomHub}
\label{fig:architecture}
\end{figure}

\section{Phạm vi sản phẩm khả dụng tối thiểu}
\label{section:mvp}

Phạm vi sản phẩm khả dụng tối thiểu là thực hiện trọn luồng mượn phòng trên một điện thoại Android: tìm phòng, gửi yêu cầu, phê duyệt, sử dụng và kết thúc. Ứng dụng có bản cài đặt chạy độc lập, còn các quy tắc nghiệp vụ chính được kiểm tra bằng ca kiểm thử tương ứng. Bảng~\ref{tab:mvp} phân biệt chức năng đã triển khai với phần cần kiểm chứng thêm trên nhiều thiết bị hoặc với khóa thật.

\begin{table}[H]
\centering
\caption{Phạm vi sản phẩm khả dụng tối thiểu}
\label{tab:mvp}
\small
\begin{tabularx}{\textwidth}{>{\raggedright\arraybackslash}p{3.4cm} X >{\centering\arraybackslash}p{1.9cm}}
\toprule
\rowcolor{RoomHubLight}
\textbf{Nhóm chức năng} & \textbf{Nội dung} & \textbf{Trạng thái} \\
\midrule
Tài khoản & Đăng nhập, đăng ký, đổi hoặc khôi phục mật khẩu, hồ sơ và quyền quản trị tài khoản. & Đã triển khai \\
Tìm phòng & Lọc theo tầng, thời gian, sức chứa, thiết bị và loại khóa; xem sơ đồ phòng theo trạng thái. & Đã triển khai \\
Đặt và duyệt & Tạo yêu cầu, kiểm tra quy tắc; phê duyệt, từ chối, đổi phòng và hủy theo thời hạn. & Đã triển khai \\
Đặt lặp hàng tuần & Tạo tối đa tám lượt, quản lý các lượt tương lai và duyệt theo chuỗi. & Đã triển khai \\
Phòng và bảo trì & Quản lý thông tin phòng, lịch bảo trì, yêu cầu báo sự cố và lịch sử sử dụng. & Đã triển khai \\
Truy cập phòng & Thỏa thuận lịch nhận khóa cơ/thẻ từ; cấp quyền và tạo mã PIN bảy chữ số cho phòng khóa số. & Đã triển khai \\
Điểm danh & Ghi nhận việc vào, ra phòng và xử lý trường hợp người đặt không đến. & Đã triển khai \\
Thông báo, cấu hình & Thông báo theo sự kiện, thông báo ưu tiên và điều chỉnh quy tắc nghiệp vụ. & Đã triển khai \\
\midrule
Đồng bộ nhiều thiết bị & Dùng dịch vụ dữ liệu chung; hiện đã kiểm tra với máy chủ và cơ sở dữ liệu trong môi trường cục bộ. & Đang hoàn thiện \\
Kiểm chứng khóa thật & Xác nhận khóa vật lý nhận và áp dụng mật khẩu theo đúng thời hạn. & Giai đoạn sau \\
Chức năng mở rộng & Thống kê mức sử dụng, gợi ý phòng nâng cao và phương án điểm danh bổ sung. & Giai đoạn sau \\
\bottomrule
\end{tabularx}
\end{table}

Quản trị viên có thể điều chỉnh khoảng thời gian được phép đặt trước. Cấu hình ban đầu cho phép đặt trước từ một đến ba ngày; nếu cho phép đặt cùng ngày, giới hạn tối thiểu có thể giảm về không. Tổng số yêu cầu đang chờ xử lý và lượt đã duyệt sắp sử dụng của mỗi người dùng không vượt quá hai. Việc hủy và đổi phòng phải tuân theo thời hạn trước giờ bắt đầu. Với khóa số, mã PIN có hiệu lực từ mười phút trước giờ bắt đầu đến mười phút sau giờ kết thúc. Nếu chưa có check-in sau thời gian chờ mặc định mười lăm phút kể từ giờ bắt đầu, quản trị viên có thể xác nhận vắng mặt. Các mốc thời gian này được kiểm tra theo cấu hình hiện hành. Vòng đời một lượt đặt phòng được minh họa ở Hình~\ref{fig:state}.

\begin{figure}[H]
\centering
\includegraphics[width=0.98\textwidth]{trang_thai_booking.pdf}
\caption{Vòng đời trạng thái của một lượt đặt phòng}
\label{fig:state}
\end{figure}

\section{Tổ chức và phân chia trách nhiệm}

Mã nguồn ứng dụng được chia thành phần hạ tầng dùng chung, các thành phần giao diện tái sử dụng và mười module nghiệp vụ. Mỗi module tập hợp màn hình, dịch vụ xử lý và mô hình dữ liệu của một phạm vi chức năng. Cách tổ chức này giúp các thành viên phụ trách những phần khác nhau mà vẫn dùng chung quy tắc về tài khoản, lưu trữ và trao đổi dữ liệu.

Khi hoạt động cục bộ, thao tác nghiệp vụ đọc và ghi dữ liệu trên điện thoại. Ở chế độ đồng bộ, ứng dụng trao đổi dữ liệu với máy chủ qua giao diện mạng và sử dụng phiên đăng nhập để xác định quyền truy cập. Hai chế độ cùng phục vụ một luồng thao tác của người dùng; việc kiểm chứng trên nhiều thiết bị vẫn là công việc tiếp theo.

\section{Các module nghiệp vụ}
\label{section:modules}

Bảng~\ref{tab:modules} trình bày trách nhiệm của mười module nghiệp vụ. Các mục tiếp theo giải thích cách chúng phối hợp trong những luồng chính của ứng dụng.

\begin{table}[H]
\centering
\caption{Phân chia module của ứng dụng}
\label{tab:modules}
\small
\begin{tabularx}{\textwidth}{>{\raggedright\arraybackslash}p{4.1cm} X}
\toprule
\rowcolor{RoomHubLight}
\textbf{Module} & \textbf{Trách nhiệm chính} \\
\midrule
Xác thực & Đăng nhập, đăng ký, đổi và khôi phục mật khẩu. \\
Quản lý tài khoản & Tìm kiếm, cập nhật, phân quyền và vô hiệu hóa tài khoản. \\
Quản lý phòng & Lưu thông tin phòng, tìm kiếm và trình bày sơ đồ theo tầng. \\
Đặt phòng & Tạo, duyệt, hủy, đổi phòng, đặt lặp và theo dõi việc sử dụng. \\
Lịch và bảo trì & Quản lý lịch bảo trì, tiếp nhận yêu cầu báo sự cố. \\
Quyền truy cập & Cấp quyền theo phòng và tạo mã PIN tạm thời. \\
Khóa thông minh & Gắn khóa vào phòng, trao đổi bản tin và tiếp nhận sự kiện từ khóa. \\
Thông báo & Gửi và phân loại thông báo theo sự kiện nghiệp vụ. \\
Hồ sơ cá nhân & Xem và cập nhật thông tin của người đang đăng nhập. \\
Cấu hình & Điều chỉnh giới hạn đặt phòng và các thời hạn nghiệp vụ. \\
\bottomrule
\end{tabularx}
\end{table}

\subsection{Tài khoản và quản lý người dùng}

Lần chạy đầu, ứng dụng chuẩn bị hai tài khoản thử nghiệm tương ứng với hai vai trò. Khi đăng ký, người dùng phải nhập tên đăng nhập, mật khẩu và xác nhận mật khẩu; thông tin cá nhân và mã khôi phục là tùy chọn. Hệ thống từ chối tên đăng nhập trùng hoặc mật khẩu không hợp lệ. Việc khôi phục mật khẩu chỉ áp dụng cho tài khoản đã thiết lập mã khôi phục.

Trước những thao tác có ảnh hưởng đến dữ liệu, ứng dụng kiểm tra tài khoản còn hoạt động và có đúng quyền. Quản trị viên không thể tự vô hiệu hóa hoặc xóa tài khoản của mình, đồng thời không thể xóa tài khoản còn lịch đặt đang hoạt động.

\subsection{Phòng và tra cứu phòng}

Dữ liệu trình diễn gồm 56 phòng thuộc tám tầng, mỗi tầng bảy phòng với sức chứa và thiết bị khác nhau. Khi thêm hoặc sửa phòng, hệ thống kiểm tra tầng, sức chứa và tên phòng để tránh dữ liệu không hợp lệ hoặc trùng lặp. Phòng còn gắn với yêu cầu đang hoạt động không được xóa.

Giảng viên có thể lọc phòng theo tầng, khung giờ, sức chứa, thiết bị và loại khóa. Kết quả được thể hiện trên sơ đồ tầng với màu trạng thái. Ở phía quản trị, danh sách phòng được thu gọn theo tám tầng; chọn một phòng sẽ mở các thao tác sửa, xóa, lên lịch bảo trì và xem lịch sử sử dụng.

\subsection{Đặt phòng, phê duyệt và sử dụng}

Khi gửi yêu cầu, ứng dụng kiểm tra phòng còn hoạt động, mục đích sử dụng đã được nhập, thời gian nằm trong khoảng được phép đặt và phiên sử dụng chưa kết thúc. Yêu cầu bị từ chối nếu phòng đang bảo trì, đã có lịch trùng, người đặt có lịch sử dụng khác cùng thời điểm hoặc đã đạt giới hạn yêu cầu. Hai lịch được xem là trùng khi cùng phòng và khoảng thời gian sử dụng của chúng có phần giao nhau; các khoảng chỉ tiếp giáp ở thời điểm kết thúc hoặc bắt đầu vẫn có thể được xét riêng.

Với yêu cầu lặp hàng tuần, hệ thống tạo tối đa tám lượt và kiểm tra cả chuỗi trước khi ghi nhận. Nếu một lượt vi phạm quy tắc, thao tác không được lưu dở dang. Quản trị viên duyệt hoặc từ chối yêu cầu sau khi kiểm tra lại lịch phòng và lịch bảo trì. Khi cần đổi phòng, phòng thay thế phải có cùng loại khóa, sức chứa không nhỏ hơn và đủ thiết bị so với phòng ban đầu.

Lịch sắp tới, lịch đã sử dụng và lịch đã hủy được trình bày theo trạng thái để tránh lặp thao tác ở nhiều nơi. Nếu đã hết thời gian chờ check-in mà vẫn chưa có bằng chứng vào phòng, quản trị viên có thể ghi nhận vắng mặt; khung giờ được giải phóng và quyền dùng mã tương ứng được thu hồi.

\subsection{Bảo trì và báo sự cố}

Quản trị viên có thể tạo lịch bảo trì cố định, nhưng lịch này không được chồng với bảo trì khác hoặc lượt đặt đã được duyệt. Nếu phòng đã được đặt, người quản lý phải xử lý việc đổi phòng trước khi chặn sử dụng. Trong thời gian sử dụng phòng theo một lượt đặt đã duyệt, người dùng có thể báo sự cố tại phòng và nêu lý do. Mỗi lượt sử dụng chỉ có một yêu cầu báo sự cố đang chờ xử lý; tầng và phòng liên quan được đánh dấu để quản trị viên nhận biết.

\subsection{Quyền truy cập và khóa thông minh}

Quyền tự tạo mã PIN được cấp theo từng cặp người dùng và phòng. Vì vậy, việc thu hồi quyền ở một phòng không làm thay đổi quyền ở phòng khác. Trước khi cấp mã, hệ thống đối chiếu lịch đã duyệt, người sở hữu lịch, loại khóa, thời hạn sử dụng và quyền tại phòng. Mỗi mã có bảy chữ số và thời hạn bao gồm mười phút trước giờ bắt đầu đến mười phút sau giờ kết thúc.

Ứng dụng lưu mã để người dùng nhìn thấy trước. Nếu đã có phiên kết nối OneIoT, lệnh tạo mật khẩu được gửi đến khóa gắn với phòng; nếu chưa kết nối, lệnh được giữ chờ và gửi khi phiên hợp lệ được mở. Cách xử lý này cho phép chuẩn bị mã trong lúc mất mạng, nhưng việc khóa thật nhận và áp dụng mã vẫn cần kiểm chứng riêng. Với phòng dùng khóa cơ hoặc thẻ từ, người dùng và quản trị viên thống nhất lịch nhận khóa; người dùng có thể ủy quyền người nhận hộ theo quy định.

Ứng dụng cũng tiếp nhận sự kiện truy cập do khóa phát ra để cập nhật việc vào và ra phòng. Thông tin xác thực OneIoT chỉ tồn tại trong phiên làm việc và được xóa khi phiên kết thúc.

\subsection{Thông báo, hồ sơ và cấu hình}

Các sự kiện như gửi yêu cầu, duyệt, từ chối, cấp mã hoặc đổi phòng đều tạo thông báo. Màu sắc thể hiện tính chất tích cực, tiêu cực, cần chú ý hoặc trung tính của từng sự kiện. Thông báo ưu tiên có thể được hiển thị nổi. Người dùng quản lý hồ sơ cá nhân; quản trị viên điều chỉnh các giới hạn và thời hạn nghiệp vụ. Trước khi lưu cấu hình, hệ thống kiểm tra các giá trị không âm, khoảng đặt trước hợp lệ và giới hạn yêu cầu tối thiểu là một.

\subsection{Dịch vụ đồng bộ dữ liệu}

Dịch vụ Java Spring Boot cung cấp giao diện dữ liệu cho tài khoản, phòng, cấu hình, bảo trì, lịch đặt, quyền tạo mã và thông báo. Máy chủ xác thực phiên, kiểm tra quy tắc nghiệp vụ và trả lỗi theo cùng một định dạng. Cơ sở dữ liệu H2 được dùng để kiểm thử cục bộ; cấu hình kết nối cơ sở dữ liệu từ xa đã được chuẩn bị. Bản trình diễn trên điện thoại dùng dữ liệu cục bộ để không phụ thuộc trạng thái máy chủ.

\section{Điều hướng giao diện}

Hình~\ref{fig:ui-flow} tóm tắt cách di chuyển giữa các chức năng theo vai trò. Giảng viên chọn khung giờ, lọc và chọn phòng, nhập mục đích rồi gửi yêu cầu ngay trên cùng một màn hình. Quản trị viên mở quản lý phòng và bảo trì theo tầng, chọn phòng rồi thực hiện thao tác tương ứng. Các chức năng khác được mở từ màn hình chính của từng vai trò. Các màn hình tương ứng được trình bày ngay sau sơ đồ điều hướng.

\begin{figure}[H]
\centering
\includegraphics[width=0.95\textwidth]{dieu_huong_giao_dien.pdf}
\caption{Luồng điều hướng giao diện theo hai vai trò}
\label{fig:ui-flow}
\end{figure}

\clearpage
\section{Giao diện ứng dụng hiện tại}
Các ảnh dưới đây được chụp trực tiếp từ phiên bản ứng dụng hiện tại cài trên điện thoại Android. Chúng minh họa màn hình theo vai trò, thao tác tìm và đặt phòng trên cùng một màn hình, cùng cách quản lý phòng theo tầng đã nêu ở Hình~\ref{fig:ui-flow}.

\begin{figure}[H]
\centering
\begin{subfigure}[t]{0.43\textwidth}
  \centering
  \includegraphics[width=\linewidth]{giao_dien_nguoi_dung_hien_tai.png}
  \caption{Màn hình người dùng}
\end{subfigure}
\hfill
\begin{subfigure}[t]{0.43\textwidth}
  \centering
  \includegraphics[width=\linewidth]{giao_dien_admin_hien_tai.png}
  \caption{Màn hình quản trị viên}
\end{subfigure}
\caption{Màn hình chính theo hai vai trò của ứng dụng hiện tại}
\label{fig:current-dashboards}
\end{figure}

\begin{figure}[H]
\centering
\begin{subfigure}[t]{0.43\textwidth}
  \centering
  \includegraphics[width=\linewidth]{tim_va_dat_phong_hien_tai.png}
  \caption{Chọn thời gian và lọc phòng}
\end{subfigure}
\hfill
\begin{subfigure}[t]{0.43\textwidth}
  \centering
  \includegraphics[width=\linewidth]{chon_phong_va_dat_hien_tai.png}
  \caption{Chọn phòng và nhập thông tin đặt}
\end{subfigure}
\caption{Tìm phòng và tạo yêu cầu trên cùng một màn hình}
\label{fig:current-booking}
\end{figure}

\begin{figure}[H]
\centering
\begin{subfigure}[t]{0.43\textwidth}
  \centering
  \includegraphics[width=\linewidth]{so_do_phong_hien_tai.png}
  \caption{Sơ đồ phòng theo tầng}
\end{subfigure}
\hfill
\begin{subfigure}[t]{0.43\textwidth}
  \centering
  \includegraphics[width=\linewidth]{quan_ly_phong_theo_tang_hien_tai.png}
  \caption{Danh sách tầng của quản trị viên}
\end{subfigure}
\caption{Hai cách xem phòng theo tầng ở ứng dụng hiện tại}
\label{fig:current-floors}
\end{figure}

\begin{figure}[H]
\centering
\begin{subfigure}[t]{0.43\textwidth}
  \centering
  \includegraphics[width=\linewidth]{phong_tang_mot_hien_tai.png}
  \caption{Phòng thuộc tầng 1}
\end{subfigure}
\hfill
\begin{subfigure}[t]{0.43\textwidth}
  \centering
  \includegraphics[width=\linewidth]{cau_hinh_hien_tai.png}
  \caption{Cấu hình quy tắc đặt phòng}
\end{subfigure}
\caption{Giao diện quản lý phòng và cấu hình nghiệp vụ}
\label{fig:current-admin}
\end{figure}

